\begin{figure}[t]
    \centering
    \includegraphics[width=\columnwidth]{%
        ../figures/paper/03_lorentzian_echo_slices.pdf}
    \caption{Experimental and simulated spectra.
    (a)--(f) Final excited-state probability $P_e$ versus detuning for
    $20~\mu\mathrm{s}$ root-Lorentzian (left) and echo-root-Lorentzian (right)
    pulses with $c=0.005$.  From top to bottom,
    $\Omega_0/2\pi=3$, $20$, and $40~\mathrm{MHz}$. Continuous lines show
    three-level simulations and markers show experiment. A common global
    population offset and small gain adjustment are applied to the measured
    traces for comparison. Residual baseline mismatches remain, but the
    spectral profiles otherwise agree closely. At high drive, the echo
    response shows a line-shape discrepancy and a deeper experimental
    resonance dip, providing a stronger signal for resonance detection.
    Partial refocusing of slowly varying, non-Markovian noise by the echo
    is a possible explanation for the enhanced dip contrast relative to the
    Markovian model, consistent with noise suppression in rotary-echo
    experiments~\cite{gustavsson_rotary_echo_2012}.
    }
    \label{fig:simulation-results}
\end{figure}
